% Article VII, additive recovery, revision 04, 10-09-2026.
% RESULTADO_RECUPERADO: S15, 11c_gravedad_holonomia_rev6.tex:16--126.
% The canonical cocycle c27=(1,18) is retained as an antecedent of route
% computation; this fragment proves its composition and representation.
% Finite translational memory and the cosmological amplitude s0 have
% different types. The cocycle alone does not evaluate a local density.

\section{Phase return and representation of accumulated memory}
\label{vii:sec:retorno-memoria-gravedad}

The chronology of enriched routes distinguishes position return,
orientation return, and phase return. The observable representation
identifies certain arrival states, while the cumulative register
preserves the sequence traversed. The geometric representation must
transport both data together.

\subsection{Reversal cocycle and change of sheet}
\label{vii:subsec:cociclo-memoria}

Let \(F_{\rm obs}\) be the chronological representation of TPK transport.
Its canonical blocks have the hierarchy
\[
\begin{array}{c|l}
27&\text{position return with orientation reversal},\\
54&\text{position and orientation return},\\
108&\text{complete return of the observable phase}.
\end{array}
\]
The register counts orientation reversals \(g(a)\) and
effective changes of sheet \(s(a)\). For a route,
\begin{equation}
 c(\gamma)=\sum_{a\subset\gamma}(g(a),s(a)),\qquad
 c(\gamma_2\gamma_1)=c(\gamma_2)+c(\gamma_1).
 \label{vii:eq:cociclo-ruta}
\end{equation}
The canonical computation for the block of length \(27\) expresses
\(c_{27}=(1,18)\). Composition preserves these increments in
the four blocks completing the observable period.
On forward-oriented routes the register is cumulative;
its extension to the route groupoid uses \(\mathbb Z^2\) and
\(c(\bar\gamma)=-c(\gamma)\).

\begin{proposition}[Memory of the complete observable turn]
\label{vii:prop:memoria-vuelta-completa}
Writing \(\mathcal K_{27}=F_{\rm obs}^{27}\), the lift
to the register, for \(x\in\mathcal O_{108}\) and \(z\in\mathbb Z^2\), satisfies
\[
 \widetilde{\mathcal K}_{27}(x,z)
   =(\mathcal K_{27}x,z+(1,18)),\qquad
 \widetilde{\mathcal K}_{27}^{\,4}(x,z)
   =(x,z+(4,72)).
\]
After \(N\) complete turns, the increment is \(N(4,72)\).
\end{proposition}

\begin{proof}
The observable projection of \(\mathcal K_{27}\) has order four.
The preceding canonical computation fixes \(c_{27}=(1,18)\) constantly
on that observable orbit. Additivity in
\eqref{vii:eq:cociclo-ruta} sums the four increments and gives
\(c_{108}=4(1,18)=(4,72)\).
Repetition applies the same concatenation law and produces
\(Nc_{108}\). Return of the first coordinate preserves the
increment of the second.
\end{proof}

\pagebreak % REV03: keep the affine cocycle subsection complete.
\subsection{Affine representation of the cocycle}
\label{vii:subsec:representacion-afin-cociclo}

The discrete realization uses an element \(R_4\) of order four in
\(\operatorname{Spin}^+(1,3)\), a unit timelike vector \(u\)
fixed by its vector action, and a unit of length \(\ell_0>0\).
The scale \(\ell_0\) retains its provenance in the dimensional section.
Define
\begin{equation}
 \rho_{\rm C}^{\rm disc}(\gamma)=
 \bigl(R_4^{c_g(\gamma)},\,\ell_0c_s(\gamma)u\bigr)
 \in\operatorname{Spin}^+(1,3)\ltimes\mathbb R^{1,3}.
 \label{vii:eq:representacion-discreta-memoria}
\end{equation}
The power acts on the vector through the associated vector
representation of the spin group.

\begin{theorem}[Composition and translational amplitude of return]
\label{vii:thm:traslacion-memoria}
The map \(\rho_{\rm C}^{\rm disc}\) preserves the identity,
concatenation, and inversion. On complete returns,
\begin{equation}
 \rho_{\rm C}^{\rm disc}(\gamma_{108})=(I,72\ell_0u),\qquad
 \rho_{\rm C}^{\rm disc}(\gamma_{108}^{\,N})=(I,72N\ell_0u).
 \label{vii:eq:traslacion-retorno-108}
\end{equation}
Inversion changes the sign of the translation.
\end{theorem}

\begin{proof}
The semidirect product is
\((R,a)(S,b)=(RS,a+Rb)\). Since \(R_4u=u\),
\[
\begin{split}
 \rho_{\rm C}^{\rm disc}(\gamma_2)
 \rho_{\rm C}^{\rm disc}(\gamma_1)
 &=
 \bigl(R_4^{c_g(\gamma_2)+c_g(\gamma_1)},
 \ell_0(c_s(\gamma_2)+c_s(\gamma_1))u\bigr)\\
 &=\rho_{\rm C}^{\rm disc}(\gamma_2\gamma_1).
\end{split}
\]
The cocycle of the identity is zero, and that of the inverse is its opposite.
For the complete turn, substitute \(c_{108}=(4,72)\) and
\(R_4^4=I\). Concatenation of \(N\) turns keeps the rotational
part equal to the identity and sums their translations.
\end{proof}

The amplitude \(72\ell_0\) is a length associated with a finite route.
A realization through a co-tetrad and connection transports it to its
affine holonomy. The local contribution to the gravitational equations
is obtained subsequently through the curvature of that connection, the
variation of the matter action, and the normalization of its density.
The representation thus preserves the exact cumulative content that
must enter that realization.
